\documentclass[11pt,a4paper]{article}
\usepackage[margin=24mm]{geometry}
\usepackage{amsmath,amssymb,amsthm}
\setlength{\emergencystretch}{2em}
\newtheorem{lemma}{Lemma}
\newtheorem{theorem}[lemma]{Theorem}
\newcommand{\gen}[1]{\langle#1\rangle}
\title{The transitive suffix in the later-cut equality case}
\author{Human Compression project: L044 continuation}
\date{}
\begin{document}\maketitle
\begin{lemma}[An anchored action on pairs]\label{pairs}
Let $v$ be a fixed-point-free involution on a set $A$ of size
$2t$, $t\ge2$. Let $Q\le\operatorname{Alt}(A)$ be transitive,
centralize $v$, and have an involutory string marking
$q_1,\ldots,q_{t-1}$. Suppose every maximal omission of the
induced marking on the $t$ pairs of $v$ is intransitive.
Suppose a point $a\in A$ is fixed by $q_2,\ldots,q_{t-1}$.
Then the pair image is the natural $S_t$ with its successive
transposition marking, anchored at the pair containing $a$.
Its kernel is nontrivial. If $t$ is even, $v\in Q$.
If $t$ is odd, $Q$ is the entire group of even permutations
centralizing $v$, namely $C_2^{\,t-1}\rtimes S_t$ in this
literal action on $A$.
\end{lemma}
\begin{proof}
Choose, for each pair-image generator, an edge crossing the
orbits of its omission. One edge per label forms a forest:
a cycle would contain an avoiding path between the endpoints
of one selected edge. There are $t-1$ edges on $t$ vertices,
so this forest is a tree $T$.

Removing its edge of label $i$ gives exactly the two pair-orbits
of that omission. Indeed all other tree edges belong to the
omission, so connect each of the two tree components, and
intransitivity prevents their joining. For a fixed label $i$,
its image preserves both sides of every cut obtained by deleting
an edge with a different label. Intersecting those cuts leaves
only one nonsingleton class: the two ends of edge $i$.
Thus the image is precisely that transposition.

Nonadjacent labels commute, so their tree edges are disjoint.
A vertex cannot have three incident edges, since three distinct
indices cannot be pairwise consecutive. Hence $T$ is a path.
The labels along it are successive indices, in one of the two
orders. The anchor pair is fixed by all labels except $1$,
so it is the endpoint incident with label $1$. The image of
$\gen{q_2,\ldots,q_{t-1}}$ is consequently its full stabilizer
$S_{t-1}$ in $S_t$.

Write $N$ for the kernel on pairs. If $N=1$, the subgroup
$\gen{q_2,\ldots,q_{t-1}}$ is the full pair stabilizer in $Q$.
It fixes $a$, so the orbit of $a$ has size at most $t$,
contrary to transitivity on $2t$ points. Therefore $N\ne1$.

The centralizer of $v$ is the permutation group on its pairs
with independent pair flips, $C_2^t\rtimes S_t$. A permutation
of the pairs is even on the underlying $2t$ points. Hence
$N$ is a nonzero $S_t$-invariant subspace of
\[
 E=\{(x_1,\ldots,x_t)\in\mathbb F_2^t:
                         x_1+\cdots+x_t=0\}.
\]
The invariance follows from the full pair image; conjugation
of flips permutes their coordinates. If $N$ contains a vector
whose support is nonempty and not all coordinates, exchange
one coordinate in its support with one outside. Adding the
vector to its image gives a vector supported on precisely
those two coordinates. Conjugating gives every such vector,
which together span $E$. Thus either $N=E$, or its only
nonzero vector is the all-ones vector. For even $t$, both
possibilities contain the all-ones flip $v$. For odd $t$,
the all-ones vector is not in $E$, so $N=E$. The order of
$Q$ is then $2^{t-1}t!$, equal to that of the full even
centralizer containing it; hence they coincide.
\end{proof}

\begin{theorem}[Exclusion of the transitive equality branch]
Retain the complete critical-head hypotheses, actual IP and all
intransitive ambient omissions. Suppose a later unique interior
cut has adjacent first shared label $g$ and crossing label
$j=g+1$. Write its two omission orbits as $C=C_g$ and $A$.
Suppose the exact reversed support equality gives
\[
 |A|=2t,\quad t=r-g-1,\quad
 \operatorname{Fix}(P_{g-1})=A\cup\{b\},
 \qquad P_k=G_{\le k},
\]
where $b\in C$ is the crossing endpoint. Then $Q=G_{>j}$
cannot be transitive on $A$.
\end{theorem}
\begin{proof}
Suppose $Q$ is transitive on $A$. The actual unique-cut lemma
makes it fix $C$ pointwise.
Its rank is $r-1-j=t-1$, and its action on $A$ is faithful
and even. The shared generator $\rho_g$ acts nontrivially
on $A$ and centralizes $Q$. A centralizing permutation of
a transitive action either has no fixed point or is the
identity; hence $v$ is a fixed-point-free involution.
The unique-crossing anchor on the $A$ side is fixed by all
$Q$ generators except its first one, by the complete acquired
crossing lemma.

Every pair-image omission in $Q$ is intransitive. For if
omitting a mark $q$ left its pair action transitive, adjoining
$v$ would make the action on $A$ transitive. In the full
ambient omission $G_q$, that joining action is supplied by
$\rho_g$; it preserves $A$ and restricts there to $v$.
Thus $A$ lies in a single $G_q$-orbit. The omitted mark $q$
is supported wholly on $A$, so adjoining it cannot join any
$G_q$-orbits. This contradicts full transitivity and the
intransitive omission hypothesis. All hypotheses of
Lemma~\ref{pairs} now hold.

If $t$ were even, the lemma would give $v\in Q$. On the other
hand the exterior image of $P_g$ on $A$ is exactly $\gen v$,
since $P_{g-1}$ fixes $A$ pointwise. Its head orbit is $C$,
and the acquired exact parity-fibre formula applies.
The involution $v$ is even when $t$ is even, so that formula
places its pure permutation on $A$ in $P_g$ as well.
Actual IP for the disjoint index sets $\{0,\ldots,g\}$ and
$\{g+2,\ldots,r-1\}$ gives $P_g\cap Q=1$, a contradiction.
Therefore $t$ is odd, necessarily $t\ge3$, and the lemma
makes $Q$ the full even centralizer of $v$.

Let $\{a,a'\}$ be the anchor pair, with $a'=a^v$, and
let $Q_1$ omit the first distinguished generator of $Q$.
Its pair image is the full $S_{t-1}$ on the other pairs,
while every one of its generators fixes $a$, and hence $a'$
by centralization. Since $t-1\ge2$, it has no fixed point
in any other pair. Thus its complete pointwise fixed set is
\[
                      \operatorname{Fix}(Q_1)=C\cup\{a,a'\}.
\]
The crossing $s=\rho_{g+1}$ centralizes $Q_1$, so preserves
this fixed set. It exchanges $a$ with $b$. It must fix $a'$:
it cannot send $a'$ into $C$, which would be a second crossing,
and cannot send it to $a$, whose image is already $b$.
Consequently $s$ also preserves $A\setminus\{a,a'\}$.

The old prefix $P_g$ preserves that complement too, because
its exterior image is $\gen v$. Since $P_g$ is transitive
on $C$, the crossing attaches $a$ to its head orbit and
$\rho_g$ then attaches $a'$. No further point can join.
Hence the actual next head orbit is exactly
\[
                         C_{g+1}=C\cup\{a,a'\}.
\]
The exterior image of $P_{g+1}$ contains the restriction
$v'=v|_{A\setminus\{a,a'\}}$, coming from $\rho_g$.
This is a nonidentity even involution: it flips $t-1\ge2$
pairs. The exact parity-fibre formula therefore puts its
pure permutation in $P_{g+1}$. But the entire even-flip kernel
belongs to $Q$, so $v'\in Q$ also. This contradicts actual
IP for the disjoint intervals $\{0,\ldots,g+1\}$ and
$\{g+2,\ldots,r-1\}$. The odd case is excluded as well.
\end{proof}

\paragraph{Attribution and scope.}
The orbit forest and cut saturation reuse the acquired actual-orbit
method, here with one edge per pair-image label. The elementary
invariant-subspace argument and the two successive pure-fibre/IP
exclusions are the
present continuation. The complete prefix and unique-crossing proofs
remain paid. No classification of projected C-groups is assumed.
This theorem applies to the equality interfaces from the separately
submitted first-later-cut reduction, but states its hypotheses
explicitly and does not depend on that packet's acceptance.
The transitive equality suffix is excluded in full. The intransitive
suffix, endpoint cuts, opposite heads and terminal singleton still
require analysis. No whole-original, question, formal or numerical credit
is claimed.
\end{document}
